\clearpage
\fancyhead[L]{\small VOYAGE DE HUMBOLDT ET BONPLAND.}
\hypertarget{native-517-lower}{}\label{native-517-lower}
\begin{center}\small Göttingische gelehrte Anzeigen. No. 91. pp. 898..904. 10 June 1809.\\[3pt]\rule{.72\linewidth}{.3pt}\end{center}

\emph{Voyage} \textsc{de Humboldt} \emph{et} \textsc{Bonpland}. \emph{Quatrième Partie, Astronomie et Magnetisme. Premier Volume, contenant un recueil d'observations astronomiques, d'opérations trigonométriques, et de mesures barométriques, faites pendant le cours d'un voyage aux régions équinoxiales du Nouveau-Continent, depuis} 1799 \emph{jusqu'en} 1803. First instalment 1808, second instalment 1809. Large quarto, 279 pages. Available from \textsc{Schöll} and \textsc{Cotta} in Paris and Tübingen.

The importance and astonishing number of the geographical positions that our celebrated countryman \textsc{von Humboldt} determined on his journey through tropical America would alone be enough to rank it among the foremost scholarly expeditions to the New World. What makes the undertaking unique is that so very much was accomplished \emph{by one man}, and that all of it was only one of the aims that this exceptional man pursued with equal enthusiasm. We could wish for nothing more than for the wealth of observations Mr \textsc{von Humboldt} brought back to Europe to be analysed and used without hur\SourceLeaf{520}ry or superficiality, with as much conscientiousness and care as possible. Fortunately, Mr \textsc{von Humboldt} found in Mr \textsc{Oltmanns} a man who was equal to this lengthy and difficult task and undertook it with exemplary fidelity and unflagging diligence. Mr \textsc{Oltmanns} had already given valuable proof of these admirable qualities. Mr \textsc{von Humboldt} could entrust this young astronomer with deriving the results from his observations without hesitation. Mr \textsc{von Humboldt}, however, went further. So that everyone could judge for themselves the quality of his observations, the accuracy of \textsc{Oltmanns}’s calculations and the relative reliability of the results, he decided to publish all his original observations in full, alongside the main steps of \textsc{Oltmanns}’s calculations. These rare observations, so important to the geography of America and secured amid such unusual hardships and dangers, certainly deserve this honour. Even so, \emph{most} of them are of such a kind that one could hardly hope to derive appreciably more accurate results from them in the future than they already yield.

According to a report by Mr \textsc{Oltmanns}, Mr \textsc{von Humboldt} himself gives the necessary information about the instruments he took on his journey in the introduction to this work. We do not know whether the introduction has already been printed. We gather only from a few incidental notes that the instruments most often used were a chronometer by \textsc{Berthoud} and a sextant by \textsc{Ramsden}. A small two-inch sextant by \textsc{Troughton} was used from time to time, and in a few places we also find mention of a small quadrant by \textsc{Bird}. The determinations of latitude are based most often on altitudes of the Sun, more rarely on altitudes of stars. The longitudes were determined mostly by chronometer, though at the more important places eclipses of Jupiter’s satellites and lunar distances were also observed, and at Cumana a solar eclipse was observed once as well. Understandably, the measured latitudes cannot claim the degree of accuracy obtainable with fixed instruments or repeating circles, and none of the longitude determinations can be fully guaranteed to within one minute of arc. This, of course, cannot lessen the value of the observations in a country where the positions of many of the places Mr \textsc{von Humboldt} located had previously been uncertain by several degrees and where most had never been determined at all.

It would be out of keeping with the purpose of our Göttingische gelehrte Anzeigen to follow our traveller step by step and give a detailed account of all his determinations. We will therefore give a brief account of the three books that make up the two instalments before us. The \emph{first book} covers the stay in Spain and the crossing to South America. Mr \textsc{von Humboldt} determined the positions of the following places: Barcellona, Montserrat, Col de Balaquet, Venta de la Sienita, Valencia, the ruins of Sagunt, Madrid, Aranjuez, Corunna and St. Croix on Teneriffa. For some of these places good determinations by other observers already exist. Their excellent agreement with Mr \textsc{von Humboldt}’s determinations shows how much confidence the latter deserve. Here, and later for every place for which other determinations were available, whether older or newer, Mr \textsc{Oltmanns} set them alongside \textsc{Humboldt}’s in order to confirm or correct them by his. These discussions are especially important and valuable for the places that must be regarded as the key points for all the longitude determinations made in South America, namely Cumana and Caraccas. They therefore make up a considerable part of the \emph{second} and \emph{third books}. After a very careful analysis of \textsc{Humboldt}’s observations for these two important places, \textsc{Oltmanns} arrives at the following results:
\SourceLeaf{521}
\begin{center} \begin{tabular}{@{}lll@{}} \textit{Cumana}, Long. & $4^{\mathrm h}\,26^{\mathrm m}\,0^{\mathrm s}$ & Lat. $10^\circ\,27'\,52''$ N. \\ \textit{Caraccas}, Long. & $4\quad 37\quad 40$ & Lat. $10\quad 30\quad 50$ N. \\ \end{tabular} \end{center}
\noindent The rest of the second book is taken up with the determinations for the islands of Tabago and Trinidad, the the mouths of the Orinoko and various places on the coast of New Andalusia. The rest of the third book concerns the great journey into the interior of South America. On that journey Mr \textsc{von Humboldt} determined the positions of a large number of remarkable places on the Orinoko and Cassiquiare and established beyond doubt that the former is connected to the Amazon River by the latter and the Rio Negro. This part of \textsc{Humboldt}’s journey is all the more important because all earlier maps of these regions were extremely inaccurate. Some even gave latitudes that were wrong by more than 4 degrees.

The value of astronomical observations from these parts of America can only be properly appreciated when one considers the many difficulties observers face there. In Caraccas the weather is so changeable that it must be counted a rare stroke of luck when an expected observation succeeds. Of the twenty-seven nights Mr \textsc{von Humboldt} patiently stayed up to observe eclipses of Jupiter’s satellites, most were lost altogether. Sometimes the sky is still completely clear 5 minutes before an expected event, and in an instant it is overcast. On 16 December 1799, when a satellite was due to be eclipsed, Mr \textsc{von Humboldt} already had the planet in his telescope’s field of view under a perfectly clear sky. His assistant \textsc{Bonpland} was already counting the seconds. Just one more minute and the observation would have been made. But suddenly a fog arises so thick that even the planet can no longer be seen. In the wilds of the Orinoko the observer’s situation is stranger still. During the night, writes Mr \textsc{von Humboldt}, we were frightened by the roaring of tigers all around us, accompanied by the howling of our dogs. The crocodiles, drawn by the fire our Indigenous guides had lit on the bank, reared halfway out of the water, as if they wanted to observe us. But in time one grows indifferent even to such scenes through repeated exposure.

Appended to the second book is a paper by Mr \textsc{von Humboldt}, entitled \emph{Essai sur les refractions dans la zone torride}. Can the same refraction table be used in different climates? He considers this interesting question from several points of view, some of them new. Mr \textsc{von Humboldt} first considers the influence of varying proportions of the constituents of the atmosphere and of the state of its humidity. He shows that they give no reason to expect refraction to differ from one climate to another. More important is how the law governing the decrease in temperature in the higher layers of air affects refraction at low altitudes and at the horizon. It is therefore of the greatest importance to ask whether this law is the same in the different zones of the Earth. Mr \textsc{von Humboldt} carried out several very remarkable experiments that directly address this question. They agree to a surprising degree on the decrease in temperature in the torrid zone. The combined result is a height of 122 toises for a decrease of 1 degree Réaumur. This is almost exactly what \textsc{Gay-Lussac} found in a balloon at a height of 6980 metres above Paris. It therefore seems that in the warm season the law is exactly the same in the temperate and torrid zones. At lower thermometer readings, however, the temperature seems to fall more slowly with height. In the absence of direct observations on this, Mr \textsc{von Humboldt} adds the results Mr \textsc{Mathieu} derived from two well-known observations made \SourceLeaf{522}in Torneå in very severe cold, using \textsc{Laplace}’s hypothesis on the decrease in air density. For a decrease of 1 degree Réaumur they give a height of $156\frac{1}{2}$ toises. This result would lead one to expect that a different table from the one for the temperate zone is needed for refraction in the frigid zone, rather than in the torrid zone. That is just the opposite of what \textsc{Bouguer} on the one hand and \textsc{Maupertuis} and \textsc{Lemonnier} on the other maintained. It is known, however, that \textsc{Bouguer} himself altered his observations to make them agree with the supposedly smaller refractions. There is also evidence to the contrary from \textsc{Legentil}’s observations in Pondichery, likewise in the torrid zone, recalculated by \textsc{Delambre}. They agree quite well with \textsc{Bradley}’s table and disagree with \textsc{Bouguer}’s. Finally, Mr \textsc{von Humboldt} himself made several observations of refraction in Cumana, Caraccas, Acapulco and on the island of Cuba, which give values consistently higher than \textsc{Bouguer}’s table. For all these reasons it seems best to keep using in the torrid zone the refraction table constructed in the temperate zone until a skilled observer, equipped with instruments appropriate to modern observing practice, has made a series of numerous observations covering the most varied states of the atmosphere in the torrid zone. In the temperate and frigid zones, on the other hand, the lower the thermometer reading in the lowest layer of air, the more slowly the temperature falls in the higher layers. This will make it necessary to change the way our refraction tables account for temperature. Alas! \textsc{Delambre}’s remark greatly discourages any hope of ever settling the matter. Even with the barometer and thermometer readings almost unchanged, refraction at the horizon often varies by 4 minutes, and no reason for this is known.

We must also note that the printing of this work by \textsc{Humboldt} has an outward elegance befitting its intrinsic worth. We wish we could say as much for the accuracy of the printed numerical data themselves, but we were sorry to find that in this respect the proofreading was rather careless. Publishing the original observations is valuable only insofar as all the figures can be trusted. It is therefore to be expected that the editor will make another careful comparison with the manuscript and conscientiously list all the misprints at the end.
